\documentclass[11pt,a4paper]{article}
\usepackage{agpstyle}

\title{\bfseries A Markov Decision Process Formulation of the\\ Continuous Art Gallery Problem\\[4pt]
\large Document 1 of 2: States, Actions, Dynamics and Reward}
\author{AGP Project --- Design Note}
\date{\today}

\begin{document}
\maketitle

\begin{abstract}
\noindent
We give a precise Markov Decision Process (MDP) formulation of the continuous Art Gallery
Problem (AGP) with holes, as the decision problem solved by the PPO agent in
\code{Polygons dataset/ppo\_agp.py}. The formulation is deliberately built so that a
\emph{staged} training curriculum --- first saturate the gallery with guards, then prune
them --- is expressible purely as a reweighting of a single fixed reward basis, with no
change to the state space, the action space or the transition kernel. Two structural
results organise the design. First, the coverage set function is monotone submodular, so
the \emph{exclusive area} of a guard (what is lost if it is removed) can only shrink as
more guards are added; redundancy is therefore a quantity that only becomes measurable
\emph{after} saturation, which is the formal justification for the saturate-then-prune
ordering. Second, both dense reward components we use are potential-based in the sense of
Ng et al., and their discounted sum telescopes \emph{exactly} into the sparse terminal
objective ``cover the gallery with as few guards as possible''. Consequently the dense
signal introduces no bias, and the curriculum controls only \emph{when} a gradient exists,
never \emph{what} is optimal. Auxiliary terms that are not potential-based are identified
explicitly, and we adopt the rule that precisely those terms must be annealed to zero.
The companion document specifies the three-stage PPO curriculum built on top of this MDP.
\end{abstract}

\tableofcontents

% ===========================================================================
\section{Scope and relation to the existing code}
% ===========================================================================

\subsection{The problem}

An \emph{art gallery} is a simple polygon, possibly with polygonal holes. A \emph{guard}
is a point inside the gallery; it sees every point connected to it by a segment lying
entirely inside the gallery. The classical AGP asks for the smallest set of guards whose
union of visibility regions is the whole gallery. Chv\'atal's theorem
\cite{chvatal1975}, with Fisk's proof \cite{fisk1978}, gives the tight worst-case bound of
$\lfloor n/3 \rfloor$ vertex guards for a simple $n$-gon, but finding a \emph{minimum}
guard set is NP-hard \cite{leelin1986}, hard to approximate below a logarithmic factor
\cite{eidenbenz2001}, and in the continuous (point-guard) setting the decision problem is
$\exists\mathbb{R}$-complete \cite{abrahamsen2018}. Exact solvers exist but scale poorly
and are usually restricted to vertex or grid guard candidates
\cite{couto2011,tozoni2016}. This combination --- geometric, combinatorial, NP-hard,
cheap to simulate, with a natural incremental decision structure --- is exactly the profile
that makes a learned construction heuristic attractive
\cite{mazyavkina2021,kool2019attention}.

\subsection{Why write the MDP down}

The current implementation already \emph{is} an MDP, but it was assembled incrementally:
reward terms were added to fix specific failure modes (aiming outside the polygon,
stacking guards, deterministic-evaluation deadlock) rather than derived from a stated
objective. Writing the MDP down serves three purposes.

\begin{enumerate}[label=(\roman*)]
  \item It separates the \emph{objective} (what we want) from the \emph{shaping} (how we
        make it learnable), which in turn tells us which reward terms may stay fixed
        forever and which must be annealed away.
  \item It makes the staged curriculum a property of the \emph{weights}, not of the
        environment, so a single environment implementation serves all stages. Every
        locomotion paper we draw on does exactly this: \citet{wang2025dribble} and
        \citet{xu2025striker} both keep one reward basis and publish one weight column per
        stage.
  \item It exposes the action space that the guard-removal idea requires. Pruning is not
        expressible in the current $\Aspace = [0,1]^2$ ``place a guard'' action space at
        all, so the extension has to be stated before it can be trained.
\end{enumerate}

\subsection{The two documents}

This document defines $\Mdp$: state, action, dynamics, observation, reward basis and
constraints, together with the calibration of the reward weights. The companion document,
\emph{A Three-Stage PPO Curriculum for the Art Gallery Problem}, fixes the weight schedule
per stage, the stage-transition criteria, the network architecture and the PPO
hyperparameters, and gives the evaluation and ablation protocol.

% ===========================================================================
\section{Geometric preliminaries}
% ===========================================================================

\subsection{Gallery instances}

\begin{definition}[Gallery]
A \emph{gallery instance} is a pair $\Gallery = (\Poly, \Holes)$ where $\Poly \subset \Real^2$
is a simple polygon with vertex list $(p_1,\dots,p_n)$ and
$\Holes = \{H_1,\dots,H_m\}$ is a finite set of pairwise disjoint simple polygons with
$H_i \subset \operatorname{int}\Poly$. The \emph{free space} is the closed region
\[
  \Free(\Gallery) \;=\; \Poly \setminus \bigcup_{i=1}^{m} \operatorname{int} H_i ,
  \qquad
  A \;=\; \Area{\Free} \;>\; 0 ,
\]
with $\Area{\cdot}$ denoting Lebesgue measure in the plane. We write $\partial\Free$ for
its boundary (outer boundary together with all hole boundaries) and
$\dbd(q) = \operatorname{dist}(q, \partial\Free)$.
\end{definition}

In the repository an instance is a JSON record with keys \code{outer\_points},
\code{holes}, \code{area} and \code{difficulty}, produced by
\code{generate\_polygon\_dataset.py} in four difficulty tiers (\code{easy}, \code{mid},
\code{hard}, \code{extreme}) that scale the vertex count ($8$--$12$ up to $24$--$40$), the
irregularity and spikiness of the outer boundary, and the number of holes ($1$--$3$ up to
$8$--$14$). The tiers will serve as the instance-difficulty curriculum in Document 2.

\subsection{Visibility and coverage}

\begin{definition}[Visibility polygon]
For $g \in \Free$, the visibility polygon is
\[
  \Vis(g) \;=\; \bigl\{\, q \in \Free \;:\; \overline{gq} \subseteq \Free \,\bigr\}.
\]
\end{definition}

\code{Security.get\_area\_coverage\_of\_guard} computes $\Vis(g)$ exactly by the standard
angular ray-sweep: rays are cast at every polygon and hole vertex angle, perturbed by
$\pm\varepsilon$ to capture the two sides of each occluding reflex vertex, the nearest
boundary hit along each ray is retained, and the hits are assembled in angular order and
clipped back to $\Free$. The result is an exact polygon, not a raster approximation ---
this matters, because it means every area quantity below is computed in continuous
geometry even though the agent \emph{perceives} the gallery as an image.

\begin{definition}[Coverage set function]
For a finite $S \subset \Free$ let
\[
  \Cov(S) \;=\; \bigcup_{g \in S} \Vis(g),
  \qquad
  f(S) \;=\; \frac{\Area{\Cov(S)}}{A} \in [0,1],
  \qquad
  \Unc(S) \;=\; \Free \setminus \Cov(S).
\]
We call $f(S)$ the \emph{coverage fraction} and $u(S) = 1 - f(S)$ the \emph{uncovered
fraction}.
\end{definition}

\subsection{Submodularity, and why pruning must come second}

\begin{proposition}[Monotone submodularity]\label{prop:submod}
$f$ is normalised ($f(\emptyset) = 0$), monotone ($S \subseteq S' \Rightarrow f(S) \le f(S')$)
and submodular: for all $S \subseteq S' \subset \Free$ and all $g \in \Free$,
\[
  f(S \cup \{g\}) - f(S) \;\ge\; f(S' \cup \{g\}) - f(S').
\]
\end{proposition}

\begin{proof}
$f$ is a (continuous) coverage function: $f(S) = \nu\bigl(\bigcup_{g \in S} \Vis(g)\bigr)$
for the finite measure $\nu(\cdot) = \Area{\cdot}/A$. Normalisation and monotonicity are
immediate. For submodularity, write the marginal gain as the measure of the part of
$\Vis(g)$ not already covered,
$f(S \cup \{g\}) - f(S) = \nu\bigl(\Vis(g) \setminus \Cov(S)\bigr)$.
Since $S \subseteq S'$ implies $\Cov(S) \subseteq \Cov(S')$, we get
$\Vis(g)\setminus\Cov(S') \subseteq \Vis(g)\setminus\Cov(S)$ and the claim follows by
monotonicity of $\nu$.
\end{proof}

Two quantities drive the whole design. When a guard is \emph{added} to the current set
$S$, the relevant number is its \emph{marginal gain}; when a guard already in $S$ is
considered for \emph{removal}, the relevant number is its \emph{exclusive area}.

\begin{definition}[Marginal gain and exclusive area]
For $g \notin S$ and $g_j \in S$,
\[
  \delta^{+}(g \mid S) \;=\; f(S \cup \{g\}) - f(S)
  \;=\; \nu\bigl(\Vis(g)\setminus\Cov(S)\bigr),
  \qquad
  \delta^{-}(g_j \mid S) \;=\; f(S) - f(S \setminus \{g_j\}).
\]
We call $\delta^{-}(g_j \mid S)$ the \emph{redundancy deficit} of $g_j$; it is exactly the
coverage that would be lost by deleting $g_j$. A guard with $\delta^{-} = 0$ is
\emph{strictly redundant}.
\end{definition}

\begin{proposition}[Redundancy is only revealed after saturation]\label{prop:prune}
Let guards be inserted in the order $g_1, g_2, \dots, g_K$ and write
$S_j = \{g_1,\dots,g_j\}$. Let $\delta^{+}_j = \delta^{+}(g_j \mid S_{j-1})$ be the gain at
insertion time, and $\delta^{-}_j = \delta^{-}(g_j \mid S_K)$ the redundancy deficit at the
end. Then for every $j$,
\[
  \delta^{-}_j \;\le\; \delta^{+}_j ,
\]
and more generally $\delta^{-}(g_j \mid S)$ is non-increasing in $S$ for $S \ni g_j$.
\end{proposition}

\begin{proof}
Put $B = S_K \setminus \{g_j\}$ and $A' = S_{j-1}$. Then $A' \subseteq B$ and, by
\cref{prop:submod},
$\delta^{+}_j = f(A' \cup \{g_j\}) - f(A') \ge f(B \cup \{g_j\}) - f(B) = \delta^{-}_j$.
The monotonicity claim is the same inequality applied to two nested sets.
\end{proof}

\Cref{prop:prune} is the formal statement of the intuition behind the staged curriculum
requested for this project. A guard that was worth placing can only become \emph{cheaper}
to delete as the gallery fills up --- never more expensive. Therefore:

\takeaway{%
\textbf{Consequence for the curriculum.} Deletion decisions taken early in an episode are
evaluated against an optimistic (too large) redundancy deficit, and deletion decisions
taken after saturation are evaluated against the smallest deficit that guard will ever
have. Spending the first stage of training purely on \emph{coverage}, and only then
switching on the delete action and the guard cost, is not merely a convenience: it is the
regime in which the pruning signal is maximally informative and least noisy.}

Two further consequences are worth recording. The greedy policy ``repeatedly add the
guard of maximum marginal gain'' inherits the $1 - 1/e$ guarantee for budgeted coverage
\cite{nemhauser1978} and the $\ln$-factor guarantee for full cover \cite{feige1998}, which
is essentially optimal for set cover; so a well-trained \stageI{} stage should approach,
not exceed, greedy. The gains available to a learned policy come from (a) choosing
\emph{continuous} guard positions rather than a finite candidate set and (b) the pruning
that greedy never performs, since greedy solutions are not minimal --- a greedy cover can
contain strictly redundant elements. This is precisely the slack that \stageII{} targets.

% ===========================================================================
\section{The decision process}
% ===========================================================================

\subsection{A contextual family of deterministic MDPs}

A single gallery induces a deterministic MDP; the learning problem is over a
\emph{distribution} of galleries. We therefore use the contextual MDP formalism
\cite{hallak2015contextual}.

\begin{definition}[AGP contextual MDP]
Let $\Dctx$ be a distribution over gallery instances. The AGP contextual MDP is
\[
  \Mdp \;=\; \bigl\langle\, \Dctx,\; \Sspace,\; \Aspace,\; \Tkernel,\; \Rfun,\; \gamma,\; \Ospace,\; \mathcal{O} \,\bigr\rangle,
\]
where for each context $\Gallery \sim \Dctx$ the tuple
$\langle \Sspace_\Gallery, \Aspace, \Tkernel_\Gallery, \Rfun_\Gallery, \gamma \rangle$ is a
deterministic finite-horizon MDP, $\Ospace$ is the observation space and
$\mathcal{O}: \Sspace \to \Ospace$ the (deterministic, lossy) observation map.
\end{definition}

All stochasticity in the environment comes from $\Dctx$ and from the policy; given the
instance and the action sequence, the trajectory is determined. This is worth stating
because it has a practical consequence: the value function's irreducible variance is
entirely \emph{instance} variance, so evaluation must always be reported per difficulty
tier, and a fixed held-out instance set is the only meaningful validation signal.

\subsection{State}

\begin{definition}[State]\label{def:state}
A state at step $t$ is
\[
  s_t \;=\; \bigl(\, \Gallery,\; \Guards_t,\; t,\; \bar{c}_t \,\bigr)
  \;\in\; \Sspace,
\]
where $\Gallery$ is the (fixed) instance, $\Guards_t = (g_1,\dots,g_{k_t}) \in \Free^{k_t}$
is the ordered tuple of currently placed guards, $t$ is the step index and
$\bar{c}_t = \max_{\tau \le t} f(\Guards_\tau)$ is the running best coverage
(used by the coverage ratchet of \cref{sec:constraints}). We write
$k_t = |\Guards_t|$, $c_t = f(\Guards_t)$, $u_t = 1 - c_t$.
\end{definition}

\begin{remark}[Why $\Guards_t$ and not only $\Cov_t$]
If the only action were $\ADD$, the covered region $\Cov_t$ would be a sufficient statistic:
the marginal gain of a new guard depends on the current coverage alone. Once $\DEL$ is
admitted this is false --- evaluating a deletion requires knowing \emph{which} guards
produced the coverage, i.e.\ the individual $\Vis(g_j)$. The guard tuple must therefore be
part of the state, and, as \cref{sec:obs} discusses, the observation must carry enough
information about it for the policy to act on deletions at all.
\end{remark}

\subsection{Action space}

\begin{definition}[Action space]\label{def:action}
\[
  \Aspace \;=\; \underbrace{\{\ADD,\ \DEL,\ \STOP\}}_{\text{mode } \mu}
  \;\times\;
  \underbrace{[0,1]^2}_{\text{point } a}.
\]
Given a state with bounding box $\operatorname{bbox}(\Free) = [x_{\min},x_{\max}] \times
[y_{\min},y_{\max}]$, the normalised point $a = (a_x,a_y)$ is decoded to
\[
  \xi(a) \;=\; \bigl(x_{\min} + a_x (x_{\max}-x_{\min}),\; y_{\min} + a_y (y_{\max}-y_{\min})\bigr).
\]
\end{definition}

The three modes are interpreted as follows.

\begin{description}[style=nextline,leftmargin=1.2em]
  \item[$\ADD$:] place a guard at $\proj(\xi(a))$, where $\proj$ is the projection into the
        free space defined in \cref{eq:proj}. A guard is \emph{always} placed.
  \item[$\DEL$:] delete the guard nearest to the decoded point,
        $j^\star = \arg\min_{j \le k_t} \lVert g_j - \xi(a) \rVert_2$. If $\Guards_t = \emptyset$
        the action is a no-op and incurs the penalty $r_{\mathrm{void}}$.
  \item[$\STOP$:] commit the current guard set and terminate the episode.
\end{description}

\begin{remark}[Deletion as a geometric pointer]
Addressing a guard by a point in the plane rather than by an index keeps the action space
fixed-size and permutation-invariant, and reuses the same continuous head as $\ADD$. This
mirrors the trick in \citet{xu2025striker}, where the velocity command is not an external
input but is \emph{derived from geometry} each step ($v^{\mathrm{target}}_t$ is the robot--ball
unit vector scaled by a magnitude). Here the same idea turns ``which of the $k_t$ guards''
--- a variable-size discrete choice --- into ``where'', a fixed-size continuous one.
\end{remark}

\begin{remark}[Backwards compatibility]
Masking the mode head to $\{\ADD\}$ recovers exactly the action space of the current
implementation, $\Aspace = [0,1]^2$. Stage~\stageI{} of the curriculum runs under this
mask, so no separate environment is needed.
\end{remark}

\paragraph{Projection into the free space.}
The bounding box strictly contains $\Free$, so a decoded point may land outside the
polygon or inside a hole. Rather than rejecting the action we project:
\begin{equation}\label{eq:proj}
  \proj(q) \;=\;
  \begin{cases}
    q, & q \in \operatorname{int}\Free, \\[3pt]
    q^\star + \theta\,(q_{\mathrm{ref}} - q^\star), & \text{otherwise,}
  \end{cases}
  \qquad
  q^\star = \arg\min_{p \in \Free}\lVert p - q\rVert_2 ,
\end{equation}
where $q_{\mathrm{ref}}$ is a fixed interior reference point (Shapely's
\code{representative\_point}) and $\theta$ is the smallest value in a fixed increasing
ladder for which the result is strictly interior. We record the indicator
$\iota_t = \Ind{\xi(a) \notin \operatorname{int}\Free}$ for the reward.

This is not cosmetic. Without projection, a deterministic (mean-action) evaluation policy
that emits an invalid point places no guard, the observation is unchanged, and the policy
re-emits the same invalid point forever --- an absorbing deadlock that produced $0\%$
coverage on at least one held-out instance before the fix. Projection guarantees the state
advances on every step, which makes the MDP \emph{progressive}: no action can leave the
state unchanged. That is a property worth preserving in any future action-space extension.

\subsection{Transition}

The transition is deterministic. With $\xi = \xi(a)$:
\[
  \Guards_{t+1} \;=\;
  \begin{cases}
    \Guards_t \oplus \proj(\xi), & \mu = \ADD, \\[2pt]
    \Guards_t \ominus g_{j^\star}, & \mu = \DEL,\ \Guards_t \neq \emptyset, \\[2pt]
    \Guards_t, & \text{otherwise,}
  \end{cases}
  \qquad
  t \mapsto t+1,
  \qquad
  \bar{c}_{t+1} = \max\bigl(\bar{c}_t,\, f(\Guards_{t+1})\bigr),
\]
where $\oplus$ and $\ominus$ are tuple append and removal. Coverage is maintained
incrementally: $\Cov_{t+1} = \Cov_t \cup \Vis(g)$ on an add, and
$\Cov_{t+1} = \bigcup_{j \ne j^\star}\Vis(g_j)$ on a delete (a re-union, which is why the
per-guard visibility polygons are cached in \code{Security.all\_guards}).

\subsection{Initial state and instance distribution}

$s_0 = (\Gallery, \emptyset, 0, 0)$ with $\Gallery \sim \Dctx$: every episode starts from
an \emph{empty} gallery. Document 2 varies $\Dctx$ across stages (difficulty tiers,
$D_4$ symmetry augmentation) and discusses a partially-saturated initialisation as an
optional accelerator for \stageII{} --- the analogue of \citet{wang2025dribble} moving the
ball spawn radius from $10$\,m to $2$\,m between stages, and of \citet{xu2025striker}
tightening the robot spawn annulus from $[0.5,4.0]$\,m to $[0.5,2.0]$\,m.

\subsection{Termination and horizon}\label{sec:term}

Let $c^{\dagger}$ be the coverage threshold at which an instance counts as solved
(currently \code{COVERAGE\_THRESHOLD}~$= 0.98$; note the design note text says $0.95$, and
the two should be reconciled). Let $K_{\max}$ be the guard budget (currently
\code{MAX\_GUARDS}~$=20$; \cref{sec:calib} recommends $24$ to leave room for
saturation) and $T_{\max}$ the step budget. An episode ends when

\begin{center}
\begin{tabularx}{\linewidth}{@{}l l X@{}}
\toprule
\textbf{Event} & \textbf{Flag} & \textbf{Condition} \\
\midrule
Solved            & \code{terminated} & $c_t \ge c^{\dagger}$ \emph{and} the agent has no pending pruning budget \\
Committed         & \code{terminated} & $\mu_t = \STOP$ \\
Guard budget spent& \code{truncated}  & $k_t \ge K_{\max}$ \\
Step budget spent & \code{truncated}  & $t \ge T_{\max}$ \\
Stalled           & \code{truncated}  & $\rho$ consecutive unproductive steps \\
\bottomrule
\end{tabularx}
\end{center}

A step is \emph{unproductive} if $\delta^{+} < \eta \cdot u_{t-1}$, i.e.\ the new guard
covers less than a fraction $\eta$ (\code{USEFUL\_AREA\_FRAC}~$= 0.10$) of what was still
uncovered; the threshold is relative to the remaining gap, so the bar falls as the gallery
fills. The patience is $\rho = $ \code{PATIENCE} $= 3$.

\begin{remark}[Solved $\ne$ terminated, once pruning exists]\label{rem:noearlystop}
In the add-only setting, terminating at $c_t \ge c^\dagger$ is correct. Once $\DEL$
exists it is not: the interesting part of the episode --- removing the guards that
saturation made redundant --- happens \emph{after} the threshold is first crossed.
In \stageII{} and \stageIII{} the episode therefore continues past first coverage and ends
on $\STOP$ or on the step budget, with $\STOP$ carrying a penalty when $c_t < c^\dagger$.
This is the reason the $\STOP$ action exists at all: it lets the agent choose the episode
length, and \cref{prop:telescope} shows a shorter episode is (weakly) preferred whenever
$\gamma < 1$.
\end{remark}

\begin{remark}[Truncation must bootstrap]\label{rem:bootstrap}
All three of the truncation conditions above end an episode at a \emph{non-absorbing}
state whose value is genuinely non-zero. Stable-Baselines3 only bootstraps the value at
such a state when the environment sets \code{info["TimeLimit.truncated"] = True}
\cite{raffin2021sb3}; the current environment does not, so every truncated episode is
treated as terminal and its return is systematically underestimated. Combined with a
negative per-guard cost this creates a real incentive to stall into the patience
truncation. This is a concrete bug to fix alongside the reward redesign.
\end{remark}

\subsection{Discount}

$\gamma$ is a design parameter, not a property of the problem. Episodes are short
($\le 40$ steps), so a discount close to $1$ is affordable and desirable: the terminal
parsimony payoff must survive back-propagation across the whole episode. We use
$\gamma = 0.99$ in \stageI{} and $\gamma = 0.997$ in \stageII{}/\stageIII{}. Note the
contrast with the locomotion papers, which run $50$\,Hz control and thousands of steps per
episode and therefore use $\gamma = 0.994$ \cite{wang2025dribble} or $\gamma = 0.995$
\cite{wang2025vision} --- effective horizons of only a few seconds. Those numbers should
\emph{not} be copied across; here a step is a guard, not $20$\,ms.

% ===========================================================================
\section{Observation: a POMDP with a privileged critic}\label{sec:obs}
% ===========================================================================

The agent does not receive $s_t$. It receives a rasterised, bounded-size view, which makes
the learning problem a POMDP \cite{kaelbling1998pomdp}. Following the asymmetric
actor--critic pattern \cite{pinto2018asymmetric} used by all three reference papers ---
\citet{wang2025dribble} give the critic simulator state and a wider first layer
($768$ vs.\ $512$); \citet{wang2025vision} and \citet{xu2025striker} both publish an explicit
actor/critic column table --- we let the critic see strictly more than the actor.

\begin{table}[t]
\centering
\small
\caption{Observation specification. The actor input is the deployable view; the critic
additionally receives privileged quantities that are exact but expensive or
non-generalisable. $H = W = 64$ (\code{IMG\_SIZE}).}
\label{tab:obs}
\begin{tabularx}{\linewidth}{@{}l c X c c@{}}
\toprule
\textbf{Component} & \textbf{Shape} & \textbf{Description} & \textbf{Actor} & \textbf{Critic} \\
\midrule
\multicolumn{5}{@{}l}{\emph{Raster channels, $H\times W$ each, \code{uint8}}}\\
\code{interior}   & $H\times W$ & $\Ind{\text{pixel} \in \Free}$ & \checkmark & \checkmark \\
\code{covered}    & $H\times W$ & $\Ind{\text{pixel} \in \Cov_t}$ & \checkmark & \checkmark \\
\code{remaining}  & $H\times W$ & $\Ind{\text{pixel} \in \Unc_t}$ & \checkmark & \checkmark \\
\code{guards}     & $H\times W$ & Gaussian splat of $\Guards_t$ & \checkmark & \checkmark \\
\code{redundancy} & $H\times W$ & at pixel of $g_j$: $255\cdot\min(1,\delta^-_j/\sigma_p)$ & \checkmark & \checkmark \\
\code{exclusive}  & $H\times W$ & region seen by exactly one guard & & \checkmark \\
\midrule
\multicolumn{5}{@{}l}{\emph{Scalars}}\\
$c_t$                     & $1$ & coverage fraction & \checkmark & \checkmark \\
$k_t / K_{\max}$          & $1$ & guard budget used & \checkmark & \checkmark \\
$t / T_{\max}$            & $1$ & step budget used & \checkmark & \checkmark \\
$\min_j \delta^-_j$       & $1$ & cheapest available deletion & \checkmark & \checkmark \\
$\bar{c}_t - c_t$         & $1$ & slack against the coverage ratchet & \checkmark & \checkmark \\
$u_t$ (exact)             & $1$ & exact uncovered fraction (not raster) & & \checkmark \\
$k_{\mathrm{ref}}$        & $1$ & greedy reference guard count & & \checkmark \\
difficulty tier           & $4$ & one-hot & & \checkmark \\
$n$, $m$                  & $2$ & vertex count, hole count & & \checkmark \\
\bottomrule
\end{tabularx}
\end{table}

\Cref{tab:obs} extends the current three-channel image in two ways that the pruning stage
requires. The \code{guards} channel is needed because the union $\Cov_t$ is
indistinguishable between one guard and five guards covering the same region --- the
existing \code{stats} vector was added for precisely this reason, but a scalar count
cannot tell the policy \emph{where} the guards are, so a $\DEL$ head has nothing to point
at. The \code{redundancy} channel makes $\delta^-_j$ directly legible: a bright guard is
expensive to delete, a dark one is nearly free. Without it, the policy would have to infer
each guard's exclusive area from the raster, which is exactly the kind of latent quantity
that \citet{wang2025vision} found had to be supplied or explicitly reconstructed --- their
ablation showed that a decoder trained \emph{after} RL rather than jointly left ball-position
predictions ``at the noise level''.

\takeaway{%
\textbf{Design rule (observability of the pruning signal).} If the reward depends on
$\delta^-_j$, the observation must expose $\delta^-_j$. Otherwise the $\DEL$ head is being
asked to solve a credit-assignment problem the perception cannot support, and the stage-II
gradient will be dominated by noise.}

The privileged $k_{\mathrm{ref}}$ deserves a note: it is the guard count produced by a
greedy max-coverage baseline on the same instance, computed once at dataset generation
time. It gives the critic an instance-difficulty normaliser --- a hard $40$-vertex gallery
with $12$ holes simply needs more guards than an easy one, and without $k_{\mathrm{ref}}$ the
value function has to learn that from the raster. It also calibrates the terminal
parsimony reward in \cref{sec:calib}. Note that $k_{\mathrm{ref}}$ must stay out of the
\emph{actor} observation, or the policy will not generalise to instances with no
precomputed baseline.

% ===========================================================================
\section{Reward}\label{sec:reward}
% ===========================================================================

\subsection{The objective we actually want}

The AGP objective is terminal and sparse:
\begin{equation}\label{eq:trueobj}
  \max_{\pi}\;
  \Ex_{\Gallery \sim \Dctx,\ \tau \sim \pi}
  \Bigl[\, w_{\mathrm{goal}}\,\Ind{c_T \ge c^{\dagger}} \;-\; w_{\mathrm{par}}\, k_T \,\Bigr].
\end{equation}
Everything else in this section exists to make \eqref{eq:trueobj} learnable without
changing what it prefers.

\subsection{Two potentials}

\begin{definition}[Coverage and parsimony potentials]
\[
  \Pot_{\mathrm{cov}}(s) \;=\; -\,w_{\mathrm{cov}}\, u(s) \;=\; -\,w_{\mathrm{cov}}\,\bigl(1 - f(\Guards(s))\bigr),
  \qquad
  \Pot_{\mathrm{par}}(s) \;=\; -\,w_{\mathrm{par}}\, k(s).
\]
The associated shaping rewards are $F_\bullet(s,a,s') = \gamma\,\Pot_\bullet(s') - \Pot_\bullet(s)$.
\end{definition}

Evaluating $F_{\mathrm{cov}}$ and $F_{\mathrm{par}}$ at $\gamma = 1$ gives a pleasant surprise:
they reproduce exactly the reward terms one would write by hand.
\[
  F_{\mathrm{cov}} =
  \begin{cases}
    +\,w_{\mathrm{cov}}\,\delta^{+}, & \mu = \ADD, \\
    -\,w_{\mathrm{cov}}\,\delta^{-}_{j^\star}, & \mu = \DEL,
  \end{cases}
  \qquad\qquad
  F_{\mathrm{par}} =
  \begin{cases}
    -\,w_{\mathrm{par}}, & \mu = \ADD, \\
    +\,w_{\mathrm{par}}, & \mu = \DEL.
  \end{cases}
\]
So ``reward new area, charge for each guard, charge for area lost on deletion, refund the
guard cost on deletion'' is not four hand-tuned terms --- it is two potentials. The
requested reward structure (``adding guards is a negative reward, removing them with no
space lost is a positive reward, with lower reward as space is lost'') falls out of
$F_{\mathrm{cov}} + F_{\mathrm{par}}$ with a single free parameter, the ratio
$w_{\mathrm{par}}/w_{\mathrm{cov}}$.

\begin{theorem}[Policy invariance, Ng et al.\ \cite{ng1999}]\label{thm:ng}
Let $\Mdp' $ be $\Mdp$ with reward $r' = r + \gamma\Pot(s') - \Pot(s)$ for any bounded
$\Pot: \Sspace \to \Real$. If $\Pot$ vanishes at absorbing states, then every optimal policy
of $\Mdp'$ is optimal in $\Mdp$ and vice versa.
\end{theorem}

\begin{proposition}[The dense reward telescopes into the sparse objective]\label{prop:telescope}
Along any trajectory of length $T$ starting from $s_0 = (\Gallery,\emptyset,0,0)$,
\[
  \sum_{t=0}^{T-1} \gamma^{t}\bigl(\gamma\Pot(s_{t+1}) - \Pot(s_t)\bigr)
  \;=\; \gamma^{T}\Pot(s_T) - \Pot(s_0).
\]
With $\Pot = \Pot_{\mathrm{cov}} + \Pot_{\mathrm{par}}$ and $\gamma = 1$ this equals
\[
  \bigl(-w_{\mathrm{cov}} u_T - w_{\mathrm{par}} k_T\bigr) - \bigl(-w_{\mathrm{cov}} \cdot 1 - 0\bigr)
  \;=\; w_{\mathrm{cov}}\, c_T \;-\; w_{\mathrm{par}}\, k_T .
\]
Hence the undiscounted episode return under the dense reward, plus the terminal bonus, is
exactly
\[
  R(\tau) \;=\; w_{\mathrm{cov}}\,c_T \;-\; w_{\mathrm{par}}\,k_T \;+\; w_{\mathrm{goal}}\Ind{c_T \ge c^{\dagger}}
  \;+\; \underbrace{\textstyle\sum_t \gamma^t r^{\mathrm{aux}}_t}_{\text{auxiliary terms}} .
\]
For $\gamma < 1$ the shaping contributes $\gamma^T \Pot(s_T) - \Pot(s_0)$ instead, which
additionally discounts the terminal potential by $\gamma^T$ and thus adds a mild preference
for short episodes --- a pressure towards fewer guards, and towards using $\STOP$ promptly.
\end{proposition}

\begin{proof}
The sum telescopes: $\sum_{t} \gamma^{t+1}\Pot(s_{t+1}) - \sum_t \gamma^t \Pot(s_t)
= \gamma^T\Pot(s_T) - \Pot(s_0)$. Substituting $u_0 = 1$, $k_0 = 0$ gives the stated value.
\end{proof}

\Cref{prop:telescope} is the central design claim of this document. The dense, per-step
signal that PPO needs is not an approximation of the sparse AGP objective; it is an exact
decomposition of it. This is the same principle \citet{wang2025vision} rely on when they
write that potential-based shaping ``guarantees that the learned behaviors remain directed
toward successful goal scoring'', and it is what lets them give a dense proxy a weight of
$500$ against a terminal goal reward of $15$ without corrupting the objective. It also
means we do \emph{not} need to anneal the coverage term away in \stageII{}, unlike
\citet{wang2025dribble}, who halve their chasing proxy ($2.0 \to 1.0$) precisely because it
is \emph{not} potential-based and does bias the optimum.

\subsection{Which terms must be annealed}

\begin{designrule}[Anneal iff non-potential]\label{rule:anneal}
A reward term may be left at a constant weight for the whole of training if and only if it
is expressible as $\gamma\Pot(s')-\Pot(s)$ for a state function $\Pot$, or it is a
component of the true objective \eqref{eq:trueobj}. Every other term biases the optimum and
must be scheduled to zero before training ends.
\end{designrule}

This gives an objective criterion for something the reference papers do by hand. In
\citet{wang2025dribble}, the reference-motion imitation reward is dropped $1.2 \to 0.0$ with
the stated reason that it is ``removed to promote exploration and discovery of effective
dribbling behaviors''; in \citet{xu2025striker}, feet-swing ($3.0 \to 0.0$), arm-swing
($1.0 \to 0.0$) and feet-yaw ($-1.0 \to 0.0$) are all zeroed at the stage boundary. All of
these are non-potential shaping terms and all of them are annealed; the terms those papers
keep alive are the ones closest to the task objective. \Cref{rule:anneal} says the pattern
is not a coincidence.

Under this rule our terms split as in \cref{tab:invariance}.

\begin{table}[t]
\centering
\small
\caption{Classification of reward terms. Only the last group needs a schedule.}
\label{tab:invariance}
\begin{tabularx}{\linewidth}{@{}l l X@{}}
\toprule
\textbf{Class} & \textbf{Terms} & \textbf{Treatment} \\
\midrule
Objective        & $r_{\mathrm{goal}}$, $r_{\mathrm{pars}}$, $r_{\mathrm{fail}}$ & fixed weights, on from \stageI{} \\
Potential-based  & $F_{\mathrm{cov}}$, $F_{\mathrm{par}}$ & fixed weights; $F_{\mathrm{par}}$ off in \stageI{} only because the $\DEL$ action is masked \\
Non-potential    & $r_{\mathrm{front}}$, $r_{\mathrm{sweep}}$, $r_{\mathrm{sep}}$ & annealed to $0$ by the end of \stageI{} \\
Constraint-like  & $r_{\mathrm{proj}}$, $r_{\mathrm{void}}$, $r_{\mathrm{stop}}$ & fixed; promoted to hard constraints in \stageIII{} \\
\bottomrule
\end{tabularx}
\end{table}

\subsection{The reward basis}

\begin{table}[t]
\centering
\small
\caption{The reward basis. One environment, one term list; the curriculum changes only the
weight column (Document 2). ``Fires'' indicates when the term is non-zero.
$\sigma_p$ is the pruning tolerance scale, $\sigma_s$ the sweep width, $d_{\min}$ the
separation radius, $\iota_t$ the projection indicator, $k_{\mathrm{ref}}$ the greedy
reference count.}
\label{tab:basis}
\renewcommand{\arraystretch}{1.35}
\begin{tabularx}{\linewidth}{@{}l l X l@{}}
\toprule
\textbf{Symbol} & \textbf{Name} & \textbf{Expression} & \textbf{Fires} \\
\midrule
\rowcolor{agpgrey}
\multicolumn{4}{@{}l}{\emph{Objective group (coverage critic)}}\\
$r_{\mathrm{cov}}$  & coverage potential & $\gamma\Pot_{\mathrm{cov}}(s') - \Pot_{\mathrm{cov}}(s)$; equals $+\delta^{+}$ on $\ADD$, $-\delta^{-}_{j^\star}$ on $\DEL$ (at $\gamma{=}1$) & every step \\
$r_{\mathrm{goal}}$ & coverage attained  & $\Ind{c_T \ge c^{\dagger}}$ & terminal \\
$r_{\mathrm{fail}}$ & coverage missed    & $-\Ind{c_T < c^{\dagger}}$ & terminal \\
\midrule
\rowcolor{agpgrey}
\multicolumn{4}{@{}l}{\emph{Parsimony group (parsimony critic)}}\\
$r_{\mathrm{par}}$   & guard-count potential & $\gamma\Pot_{\mathrm{par}}(s') - \Pot_{\mathrm{par}}(s)$; equals $-1$ on $\ADD$, $+1$ on $\DEL$ & every step \\
$r_{\mathrm{prune}}$ & clean-prune bonus     & $\exp\!\bigl(-\delta^{-}_{j^\star}/\sigma_p\bigr) \in (0,1]$ & $\DEL$ \\
$r_{\mathrm{pars}}$  & terminal parsimony    & $\bigl(k_{\mathrm{ref}} - k_T\bigr)\,\Ind{c_T \ge c^{\dagger}}$ & terminal \\
\midrule
\rowcolor{agpgrey}
\multicolumn{4}{@{}l}{\emph{Scaffolding group --- annealed to zero (\cref{rule:anneal})}}\\
$r_{\mathrm{front}}$ & frontier bonus       & $\Ind{\proj(\xi) \in \Unc_t}$ & $\ADD$ \\
$r_{\mathrm{sweep}}$ & boundary-inward prior& $\exp\!\bigl(-(\dbd(\proj(\xi)) - \mu_{\mathrm{ring}}(t))^2 / 2\sigma_s^2\bigr)$ & $\ADD$ \\
$r_{\mathrm{sep}}$   & separation penalty   & $-\Ind{\min_j \lVert \proj(\xi) - g_j\rVert < d_{\min}}$ & $\ADD$ \\
\midrule
\rowcolor{agpgrey}
\multicolumn{4}{@{}l}{\emph{Regularity group --- constraints in \stageIII{}}}\\
$r_{\mathrm{proj}}$  & aimed outside $\Free$ & $-\iota_t$ & $\ADD$ \\
$r_{\mathrm{void}}$  & deleted from empty set& $-\Ind{\Guards_t = \emptyset}$ & $\DEL$ \\
$r_{\mathrm{stop}}$  & premature commit      & $-\Ind{c_t < c^{\dagger}}$ & $\STOP$ \\
\bottomrule
\end{tabularx}
\end{table}

\Cref{tab:basis} lists every term. Three remarks on the less obvious entries.

\paragraph{$r_{\mathrm{prune}}$ and the ``lower reward as space is lost'' requirement.}
The bounded exponential kernel $\exp(-\delta^{-}/\sigma_p)$ equals $1$ for a strictly
redundant guard and decays smoothly to $0$ as the deletion becomes expensive. Bounding
each shaping term in $(0,1]$ is the single most consistent stylistic choice across the
reference papers --- \citet{wang2025dribble} write \emph{every} tracking reward as
$\exp\{-\alpha\lVert x - x^\star\rVert^2\}$, and \citet{xu2025striker} use
$\exp(-(v^{\mathrm{target}}-v)^2/\sigma)$. The practical payoff is that a weight then \emph{is}
the term's maximum per-step contribution, which makes hand-balancing a dozen terms
tractable. Note that $r_{\mathrm{prune}}$ is technically non-potential and therefore does bias
the optimum; it survives \cref{rule:anneal} only because the bias is negligible
($\sigma_p \ll 1$, so it is effectively a delta function at $\delta^- \approx 0$, i.e.\ a
tie-breaker among genuinely free deletions). If purity is preferred, it can be annealed
with the scaffolding group.

\paragraph{$r_{\mathrm{sweep}}$ and the ``edges first, then inwards'' idea.}
The ring target
\begin{equation}\label{eq:ring}
  \mu_{\mathrm{ring}}(t) \;=\; \mu_0 + (\mu_1 - \mu_0)\,\frac{t}{K_{\max}},
  \qquad 0 \le \mu_0 < \mu_1 ,
\end{equation}
sweeps the preferred distance-to-boundary outward-to-inward \emph{within} an episode: early
guards are rewarded for sitting near $\partial\Free$ (where reflex vertices and hole
boundaries create the occlusions), later ones for moving into the interior. This is a prior
on the \emph{order} of placement, not on the final configuration, and it is the AGP analogue
of the gait-phase indicators (\code{feet swing} $\times$ \code{swing period}) that
\citet{xu2025striker} use to install a walking rhythm and then delete. Both $\mu_0,\mu_1$
should be expressed as fractions of the bounding-box diagonal so the term is scale-free.
An alternative formulation replaces $\dbd$ by distance to the nearest reflex vertex,
$\min_{v \in \Reflex}\lVert \cdot - v\rVert$, which is better motivated by the geometry
(reflex vertices are exactly the sources of occlusion, and are what makes $\lfloor n/3\rfloor$
tight \cite{chvatal1975,orourke1987}) but requires computing $\Reflex$ per instance.

\paragraph{$r_{\mathrm{pars}}$ and calibration against a reference solver.}
Rewarding $-k_T$ directly makes the terminal reward's scale instance-dependent: an easy
gallery needs $3$ guards, an extreme one $15$, so the same $-k_T$ means very different
things. Measuring against the greedy baseline, $(k_{\mathrm{ref}} - k_T)$, normalises this and
gives a directly interpretable quantity --- ``guards saved relative to greedy''. This
mirrors \citet{xu2025striker}, who set their constraint threshold to
``$1.5\times$ per-step average regularization cost of a walking gait'', i.e.\ calibrated
against a known-good reference behaviour rather than in absolute units.

\subsection{Weight calibration}\label{sec:calib}

Three inequalities determine the weights; everything else is tuning.

\paragraph{(C1) The guard--area exchange rate.}
The ratio
\begin{equation}\label{eq:exchange}
  \epsilon \;=\; \frac{w_{\mathrm{par}}}{w_{\mathrm{cov}}}
\end{equation}
is the fraction of gallery area that one guard is worth. It is the single most
interpretable hyperparameter in the whole design, and it should be set from the problem,
not by sweeping: if a typical solved instance uses $k_{\mathrm{typ}}$ guards, then
$\epsilon \approx 1/(2 k_{\mathrm{typ}})$ says ``a guard is worth half of its fair share of
the gallery'', which keeps coverage dominant while still making redundancy expensive.
With $k_{\mathrm{typ}} \approx 8$--$12$ across the four tiers we take $\epsilon = 0.02$,
i.e.\ $w_{\mathrm{cov}} = 1.0$, $w_{\mathrm{par}} = 0.02$.

\paragraph{(C2) Feasibility must dominate parsimony.}
The agent must never prefer a smaller infeasible solution to a larger feasible one. Since
at most $K_{\max}$ guards can be placed, it suffices that
\begin{equation}\label{eq:dominance}
  w_{\mathrm{goal}} \;+\; w_{\mathrm{fail}} \;>\; w_{\mathrm{par}}\,K_{\max}
  \;+\; w_{\mathrm{pars}}\,k_{\mathrm{ref}}^{\max}.
\end{equation}
With $w_{\mathrm{par}} = 0.02$, $K_{\max} = 24$, $w_{\mathrm{pars}} = 0.25$ and
$k_{\mathrm{ref}}^{\max} \approx 18$ the right-hand side is $\approx 4.98$, so
$w_{\mathrm{goal}} = 4.0$ and $w_{\mathrm{fail}} = 2.0$ satisfy \eqref{eq:dominance} with margin.
This is the AGP analogue of \citet{wang2025vision} setting their fall penalty to $-1000$
against task rewards of order $10$--$500$: the hard-failure term is deliberately off-scale.

\paragraph{(C3) The pruning tolerance.}
A deletion is immediately profitable when
\[
  w_{\mathrm{par}} + w_{\mathrm{prune}}\exp\!\bigl(-\delta^{-}/\sigma_p\bigr)
  \;>\; w_{\mathrm{cov}}\,\delta^{-}.
\]
The left side is decreasing in $\delta^-$ and the right side increasing, so there is a
unique root $\epsilon^{\star}$ --- the \emph{pruning tolerance}, the largest coverage loss
the agent will accept in exchange for one fewer guard:
\begin{equation}\label{eq:tolerance}
  \epsilon^{\star}
  \;=\; \frac{w_{\mathrm{par}}}{w_{\mathrm{cov}}}
      + \frac{w_{\mathrm{prune}}}{w_{\mathrm{cov}}}\,e^{-\epsilon^{\star}/\sigma_p}
  \;\approx\; \epsilon \;+\; \frac{w_{\mathrm{prune}}}{w_{\mathrm{cov}}}\,e^{-\epsilon/\sigma_p}
  \quad\text{when } \sigma_p \ll \epsilon .
\end{equation}
With $w_{\mathrm{cov}} = 1.0$, $w_{\mathrm{par}} = 0.02$, $w_{\mathrm{prune}} = 0.15$,
$\sigma_p = 0.005$, the root is $\epsilon^{\star} \approx 0.022$: the agent will trade away
up to about $2.2\%$ of the gallery to remove a guard, and essentially nothing beyond that.
Note that $\epsilon^\star$ is dominated by the exchange rate $\epsilon$, which is the
intended behaviour --- the exponential bonus only breaks ties among nearly-free deletions.
The hard feasibility floor is enforced separately by the ratchet of
\cref{sec:constraints}, so $\epsilon^\star$ governs \emph{intermediate} states, not the
final answer.

\paragraph{Default values.}
\begin{center}
\small
\begin{tabular}{@{}l l l@{}}
\toprule
\textbf{Symbol} & \textbf{Value} & \textbf{Set by} \\
\midrule
$w_{\mathrm{cov}}$    & $1.0$   & normalisation (defines the unit of reward) \\
$w_{\mathrm{par}}$    & $0.02$  & (C1), exchange rate $\epsilon = 1/(2k_{\mathrm{typ}})$ \\
$w_{\mathrm{goal}}$   & $4.0$   & (C2) \\
$w_{\mathrm{fail}}$   & $2.0$   & (C2) \\
$w_{\mathrm{pars}}$   & $0.25$  & (C2), scale of ``guards saved vs.\ greedy'' \\
$w_{\mathrm{prune}}$  & $0.15$  & (C3) \\
$\sigma_p$            & $0.005$ & (C3), $0.5\%$ of gallery area counts as ``no space lost'' \\
$\epsilon^{\star}$    & $0.022$ & derived from (C1)+(C3) \\
$c^{\dagger}$         & $0.98$  & task specification \\
$K_{\max}$            & $24$    & budget; must exceed $k_{\mathrm{ref}}^{\max}$ with slack for saturation \\
\bottomrule
\end{tabular}
\end{center}

\subsection{Reward groups and the multi-critic split}

The rows of \cref{tab:basis} are grouped so that each group can be given its own value
head, following \citet{wang2025vision}, who report that ``using a single critic to estimate
combined rewards can cause negative interference between distinct reward components,
leading to reduced learning stability and performance'' and combine advantages as
$\Adv_{\mathrm{total}} = w_{\mathrm{goal}}\Adv_{\mathrm{goal}} + w_{\mathrm{aux}}\Adv_{\mathrm{aux}}$
with weights $2$ and $1$. Our split is
\begin{equation}\label{eq:multicritic}
  \Adv_{\mathrm{total}}
  \;=\; \beta_{\mathrm{cov}}\,\Adv_{\mathrm{cov}}
  \;+\; \beta_{\mathrm{par}}\,\Adv_{\mathrm{par}}
  \;+\; \beta_{\mathrm{reg}}\,\Adv_{\mathrm{reg}} ,
\end{equation}
with the coverage group, the parsimony group, and the scaffolding-plus-regularity group
respectively. The reason this matters here is scale: $r_{\mathrm{cov}}$ has per-step
magnitude up to $\sim 0.5$ early in an episode and $\sim 10^{-3}$ late, while
$r_{\mathrm{par}}$ is a constant $\pm 0.02$. A single critic must fit both regimes with one
output head; two heads need not. The stage-dependent $\beta$ schedule is set in Document 2.

% ===========================================================================
\section{Constraints}\label{sec:constraints}
% ===========================================================================

Two requirements are better expressed as constraints than as reward terms, giving a
constrained MDP \cite{altman1999}.

\paragraph{Validity.} $g \in \Free$ for every placed guard. This is enforced structurally
by $\proj$ (\cref{eq:proj}), with $r_{\mathrm{proj}}$ retained as a soft penalty so the policy
still learns to aim inside rather than relying on the projection.

\paragraph{The coverage ratchet.} Once the agent has achieved coverage $\bar{c}_t$ it
should not be allowed to fall back arbitrarily far while pruning:
\begin{equation}\label{eq:ratchet}
  c_t \;\ge\; \min\bigl(\bar{c}_t,\ c^{\dagger}\bigr) - \delta_{\mathrm{slack}},
  \qquad \delta_{\mathrm{slack}} = 0.01 .
\end{equation}
In \stageII{} this is implemented as an \emph{action mask}: a $\DEL$ whose
$\delta^{-}_{j^\star}$ would violate \eqref{eq:ratchet} is rejected and charged
$r_{\mathrm{void}}$. In \stageIII{} it is promoted to a state-wise constraint solved with a
constrained policy-gradient method:
\begin{equation}\label{eq:statewise}
  \max_\pi\; \Ex\Bigl[\textstyle\sum_t \gamma^t r^{\mathrm{task}}_t\Bigr]
  \quad\text{s.t.}\quad
  c^{\mathrm{cost}}(s_t, a_t) \le h \quad \forall (s_t,a_t) \in \tau,\ \tau \sim p_\pi .
\end{equation}
\Cref{eq:statewise} is the formulation \citet{xu2025striker} adopt (their Eqs.~4--5, solved
with N-P3O \cite{lee2023constrained}), and their motivation transfers directly. They
diagnose \emph{inhomogeneous credit assignment}: task reward is concentrated at one
moment of the episode (the kick), so ``a fixed regularization coefficient can impose
disproportionately large penalties at certain steps --- particularly those immediately
preceding the kick, where the agent expects high immediate future reward''. Their ablation
is the strongest part of the argument: raising the fixed coefficient from $1.0$ to $1.5$
\emph{lowered} success from $64.8\%$ to $43.4\%$, while the constrained formulation reached
$79.5\%$. AGP has exactly the same profile --- the first few guards each capture a large
share of the area and the last ones capture slivers, so any fixed per-step penalty is
mis-scaled at one end of the episode or the other.

% ===========================================================================
\section{The environment step, in full}
% ===========================================================================

\begin{algorithm}[h]
\small
\caption{One environment step of $\Mdp$}
\begin{algorithmic}[1]
\Require state $s_t = (\Gallery, \Guards_t, t, \bar c_t)$, action $(\mu, a)$, stage weight vector $w$
\State $\xi \gets \xi(a)$ \Comment{decode into gallery coordinates}
\If{$\mu = \STOP$}
  \State \Return $s_t$, $r = w_{\mathrm{stop}}\,r_{\mathrm{stop}} + \textsc{Terminal}(s_t)$, \code{terminated}
\ElsIf{$\mu = \ADD$}
  \State $\iota \gets \Ind{\xi \notin \operatorname{int}\Free}$; \; $g \gets \proj(\xi)$
  \State $\delta^{+} \gets \nu(\Vis(g) \setminus \Cov_t)$ \Comment{exact, via Shapely difference}
  \State $\Guards_{t+1} \gets \Guards_t \oplus g$; \; cache $\Vis(g)$
  \State $r \gets w_{\mathrm{cov}}\delta^{+} - w_{\mathrm{par}} + w_{\mathrm{front}} r_{\mathrm{front}} + w_{\mathrm{sweep}} r_{\mathrm{sweep}} + w_{\mathrm{sep}} r_{\mathrm{sep}} - w_{\mathrm{proj}}\iota$
\Else \Comment{$\mu = \DEL$}
  \If{$\Guards_t = \emptyset$} \State \Return $s_t$, $r = -w_{\mathrm{void}}$, continue \EndIf
  \State $j^\star \gets \arg\min_j \lVert g_j - \xi\rVert_2$; \; $\delta^{-} \gets f(\Guards_t) - f(\Guards_t \ominus g_{j^\star})$
  \If{$f(\Guards_t) - \delta^{-} < \min(\bar c_t, c^{\dagger}) - \delta_{\mathrm{slack}}$}
     \State \Return $s_t$, $r = -w_{\mathrm{void}}$, continue \Comment{ratchet violation, masked}
  \EndIf
  \State $\Guards_{t+1} \gets \Guards_t \ominus g_{j^\star}$
  \State $r \gets -w_{\mathrm{cov}}\delta^{-} + w_{\mathrm{par}} + w_{\mathrm{prune}}\exp(-\delta^{-}/\sigma_p)$
\EndIf
\State $c_{t+1} \gets f(\Guards_{t+1})$; \; $\bar c_{t+1} \gets \max(\bar c_t, c_{t+1})$
\State evaluate termination (\cref{sec:term}); on terminal add $\textsc{Terminal}(s_{t+1})$
\State \textbf{on truncation:} set \code{info["TimeLimit.truncated"] = True} \Comment{\cref{rem:bootstrap}}
\State \Return $s_{t+1}$, $r$, flags
\Statex
\Function{Terminal}{$s$}
  \State \Return $w_{\mathrm{goal}}\Ind{c \ge c^{\dagger}} - w_{\mathrm{fail}}\Ind{c < c^{\dagger}} + w_{\mathrm{pars}}(k_{\mathrm{ref}} - k)\Ind{c \ge c^{\dagger}}$
\EndFunction
\end{algorithmic}
\end{algorithm}

\begin{figure}[t]
\centering
\begin{tikzpicture}[scale=0.92]
  \def\gal{(0,0) -- (4.2,0) -- (4.2,1.4) -- (2.7,1.4) -- (2.7,2.5) -- (4.2,2.5) -- (4.2,4.0) -- (1.6,4.0) -- (1.6,3.0) -- (0,3.0) -- cycle}
  \def\hole{(0.7,0.7) rectangle (1.7,1.7)}

  % ---- panel (a)
  \begin{scope}[xshift=0cm]
    \fill[even odd rule, agpgrey] \gal \hole;
    \draw[thick, agpblue] \gal;
    \draw[thick, agpblue] \hole;
    \node[below] at (2.1,-0.35) {\small (a) $s_0$: empty gallery};
  \end{scope}

  % ---- panel (b)
  \begin{scope}[xshift=5.4cm]
    \fill[even odd rule, agpgrey] \gal \hole;
    \draw[thick, agpblue] \gal;
    \draw[thick, agpblue] \hole;
    \foreach \p in {(0.3,0.3),(0.3,2.7),(1.9,3.7),(3.9,3.7),(3.9,2.7),(3.9,0.3),(2.4,0.3),(2.1,2.1),(1.1,2.4),(3.0,1.7)}
      \fill[agpaccent] \p circle (2.6pt);
    \node[below] at (2.1,-0.35) {\small (b) \stageI{}: saturate, $k=10$};
  \end{scope}

  % ---- panel (c)
  \begin{scope}[xshift=10.8cm]
    \fill[even odd rule, agpgrey] \gal \hole;
    \draw[thick, agpblue] \gal;
    \draw[thick, agpblue] \hole;
    \foreach \p in {(0.3,0.3),(0.3,2.7),(1.9,3.7),(3.9,0.3),(3.0,1.7)}
      \fill[agpgreen] \p circle (2.6pt);
    \foreach \p in {(0.3,2.7),(3.9,3.7),(3.9,2.7),(2.4,0.3),(2.1,2.1),(1.1,2.4)}
      \node[agpaccent, opacity=0.45] at \p {$\times$};
    \node[below] at (2.1,-0.35) {\small (c) \stageII{}: prune, $k=5$};
  \end{scope}
\end{tikzpicture}
\caption{The saturate-then-prune trajectory that the MDP is built to support.
\Cref{prop:prune} says each guard's redundancy deficit $\delta^-_j$ is smallest in panel
(b), so that is the state in which deletions (crosses in panel (c)) are cheapest and the
pruning signal is least ambiguous.}
\label{fig:pipeline}
\end{figure}

% ===========================================================================
\section{Mapping onto the current implementation}
% ===========================================================================

\Cref{tab:mapping} places every reward term currently in \code{ppo\_agp.py} into the
formalism above.

\begin{table}[h]
\centering
\small
\caption{The existing reward, classified. ``Invariant'' means the term does not change the
optimal policy and may therefore keep a fixed weight (\cref{rule:anneal}).}
\label{tab:mapping}
\begin{tabularx}{\linewidth}{@{}l l c X@{}}
\toprule
\textbf{Code} & \textbf{Formal term} & \textbf{Invariant} & \textbf{Comment} \\
\midrule
\code{area\_term}      & $F_{\mathrm{cov}}$ on $\ADD$ & yes & already exactly $\delta^{+}$ \\
\code{cost\_term}      & $F_{\mathrm{par}}$ on $\ADD$ & yes & $w_{\mathrm{par}} = $ \code{GUARD\_COST} $= 0.05$, i.e.\ $\epsilon = 0.05$ --- generous; (C1) suggests $0.02$ \\
\code{placement\_term} & $r_{\mathrm{front}}$ & no & biases towards placing in uncovered spots even when a covered spot would see more; anneal \\
\code{spacing\_term}   & $r_{\mathrm{sep}}$ & no & a fixed $15\%$-of-diagonal exclusion is wrong for corridors; anneal \\
\code{invalid\_term}   & $r_{\mathrm{proj}}$ & (constraint) & keep \\
--- & $r_{\mathrm{goal}}$, $r_{\mathrm{fail}}$, $r_{\mathrm{pars}}$ & --- & \textbf{absent}: no terminal reward of any kind \\
--- & $r_{\mathrm{prune}}$, $\DEL$ & --- & \textbf{absent}: the action space has no removal \\
\bottomrule
\end{tabularx}
\end{table}

Four gaps follow, in rough order of impact.

\begin{enumerate}[label=\textbf{G\arabic*.},leftmargin=2.6em]
  \item \textbf{No terminal reward.} \Cref{prop:telescope} shows the potential terms already
        encode $w_{\mathrm{cov}}c_T - w_{\mathrm{par}}k_T$, but there is no
        $\Ind{c_T \ge c^{\dagger}}$ term, so crossing the threshold is worth nothing beyond
        the area itself. Combined with \textbf{G2} this makes giving up cheap.
  \item \textbf{Truncation is treated as termination} (\cref{rem:bootstrap}). The patience
        stall, the guard-budget exhaustion and the step limit all cut the episode at a
        state with real remaining value, and none of them sets
        \code{TimeLimit.truncated}. The value target is therefore biased low exactly on the
        trajectories the agent should learn to avoid. This is a plausible contributor to the
        training instability recorded in the design note ($\text{approx\_kl}$ in the tens,
        $\text{clip\_fraction}\approx 0.8$, collapse after $150$--$200$k steps) --- not the
        only one, but a cheap fix to rule out.
  \item \textbf{No removal action}, so the pruning stage cannot be expressed. This is the
        action-space extension of \cref{def:action}.
  \item \textbf{The observation cannot support pruning} even if the action existed: guards
        are not visible as individual objects and $\delta^{-}_j$ is not exposed
        (\cref{tab:obs}).
  \end{enumerate}

A fifth, smaller item: \code{COVERAGE\_THRESHOLD} is $0.98$ in code and $0.95$ in the prose
of \code{docs/agp\_ppo\_design.md}; the calibration in \cref{sec:calib} assumes $0.98$.

% ===========================================================================
\section{Evaluation metrics}
% ===========================================================================

The MDP return is a training signal, not a report. Evaluation on the held-out set should
report, per difficulty tier:

\begin{itemize}
  \item \textbf{Feasibility rate} $\Prob[c_T \ge c^{\dagger}]$ --- the primary number.
  \item \textbf{Mean coverage} $\Ex[c_T]$, including failures.
  \item \textbf{Guard count} $\Ex[k_T \mid c_T \ge c^{\dagger}]$, conditioned on success so that
        failures cannot flatter the count.
  \item \textbf{Competitive ratio} $\Ex[k_T / k_{\mathrm{ref}}]$ against greedy, and against an
        exact IP solution \cite{couto2011,tozoni2016} on the \code{easy} tier where one is
        computable.
  \item \textbf{Pruning yield} $\Ex[k_{\mathrm{peak}} - k_T]$: how many guards \stageII{} removes.
        This is the quantity the whole second stage exists to increase, and it is zero by
        construction for any add-only policy.
  \item \textbf{Redundancy at termination} $\Ex[\#\{j : \delta^{-}_j < \sigma_p\}]$: strictly
        redundant guards left behind. A well-pruned solution should have none.
\end{itemize}

Reporting $k_T$ without conditioning on feasibility is the easiest way to be fooled here: a
policy that stops after two guards has an excellent guard count.

% ===========================================================================
\section{Summary}
% ===========================================================================

\begin{enumerate}[label=\textbf{\arabic*.}]
  \item The AGP is modelled as a deterministic contextual MDP over gallery instances, with
        state $(\Gallery, \Guards_t, t, \bar c_t)$ and action space
        $\{\ADD,\DEL,\STOP\}\times[0,1]^2$, where $\DEL$ addresses a guard geometrically
        rather than by index.
  \item Coverage is a monotone submodular set function, so a guard's redundancy deficit can
        only shrink as the gallery fills (\cref{prop:prune}). Saturating first and pruning
        second is therefore the ordering in which the pruning signal is strongest.
  \item The dense reward consists of exactly two potentials, on uncovered area and on guard
        count. Their discounted sum telescopes into the true sparse objective
        (\cref{prop:telescope}), so the dense signal is unbiased by construction and the
        curriculum can reweight freely.
  \item A term may keep a fixed weight iff it is potential-based or part of the objective
        (\cref{rule:anneal}); this yields an explicit, principled list of which terms the
        curriculum must anneal to zero, rather than the hand-chosen schedules used in the
        reference papers.
  \item Weights are fixed by three inequalities: the guard--area exchange rate
        $\epsilon = w_{\mathrm{par}}/w_{\mathrm{cov}}$, feasibility dominance, and the pruning
        tolerance $\epsilon^{\star}$, which comes out at about $2.2\%$ of gallery area for
        the defaults.
  \item The current implementation is the $\ADD$-only special case with no terminal reward,
        two biasing shaping terms, and truncation handled as termination. All four gaps are
        addressable without changing the environment's external interface.
\end{enumerate}

\bibliographystyle{plainnat}
\renewcommand{\refname}{References}
\bibliography{refs}

\end{document}
